%---------------------------------------------------------------
\chapter{\kesimpulan}
\label{bab:6}
%---------------------------------------------------------------

In this chapter, we present the conclusions drawn from our research and recommendations for future work.


%---------------------------------------------------------------
\section{Conclusion}
\label{sec:kesimpulan}
%---------------------------------------------------------------
% Summarize answers to the three research questions:

% RQ1 (Domain Gap):
% Moral reasoning CNFs are structurally underconstrained compared to CLUTRR/ProntoQA:
% few variables, abstract predicate semantics, and the open-textured nature of moral
% norms prevent the backbone from narrowing sufficiently for formal entailment.

% RQ2 (Pipeline Bottleneck):
% LLM-generated clauses fail to achieve SAT-based entailment due to three bottlenecks:
% (1) prompt bias toward known norm strings rather than bridging concepts,
% (2) generated predicates that do not map to existing SAT variables, and
% (3) the CoT confidence threshold firing before rule injection has a meaningful effect.

% RQ3 (CoT Compensation):
% Despite SAT failure, backbone-grounded CoT still outperforms the vanilla baseline.
% The backbone facts and the dataset's explanation field provide sufficient grounding
% for the LLM to improve moral norm classification accuracy.


%---------------------------------------------------------------
\section{Future Work}
\label{sec:saran}
%---------------------------------------------------------------
% Suggestions for future research:
% 1. Richer CNF encodings: include more fine-grained moral predicates to give
%    the backbone more room to narrow toward entailment.
% 2. Structured abduction: instead of fill-in-the-blank, use a dedicated
%    retrieval step to find known moral rules from a curated knowledge base.
% 3. Larger and more capable LLMs: test whether a larger model reduces
%    the prompt bias and variable grounding failure rates.
% 4. Extending to other moral datasets: e.g., ETHICS, MoralStories, or Delphi,
%    to test generalizability of the domain gap finding.
% 5. Dynamic CoT threshold scheduling: explore adaptive γ annealing strategies
%    that give rule injection more iterations before falling back to CoT.
